% GENERATED FILE. Edit canonical lessons, then run npm run generate:books.
% canonical-source-hash: fnv1a64:fdf43fda2d0f1da1
% canonical-lessons: JA-C315-suruto, JA-C315-tokoroga, JA-C315-shikashi, JA-C315-tokorodokoro, JA-C315-miharu

\chapter{And Then}
\label{ch:ja-and-then}
\hlchaptermodality{315}

\begin{chapteropening}
\textbf{By the end of this chapter:} \emph{I can link the steps of a story with then, but and however.}

Complete the last lesson of chapter 315: I can link the steps of a story with then, but and however.
\end{chapteropening}

\section[\texorpdfstring{\ja{すると} (suruto)}{すると (suruto)}]{\ja{すると} (suruto) — thereupon; at that}

\label{lesson:JA-C315-suruto}



\begin{quote}
Before the new one: say the Japanese for to decrease, then the Japanese for to think of.
\end{quote}

\subsection*{You'll want to know: \ja{すると}}
\textbf{\ja{すると}} — \emph{suruto} — \textquotedblleft{}thereupon; at that\textquotedblright{}.

\begin{quote}
\textbf{\ja{すると、いぬが} \ja{ほえました。}} — \emph{suruto, inu ga hoemashita} — \textquotedblleft{}At that, the dog began to bark.\textquotedblright{}
\end{quote}

\begin{cousinweb}
\textbf{\ja{すると}} (\emph{suruto}) is a native Japanese word (wago), inherited rather than borrowed from Chinese, and it is usually written in kana (JMdict).
\end{cousinweb}

\begin{grammarlens}[title={What you've built}]
One more word for telling what happened.
\end{grammarlens}

\subsection*{Your turn}
\begin{itemize}
\raggedright
  \item \emph{Say it:} \emph{suruto}
  \item \emph{Say it:} \emph{suruto}, once more
  \item \emph{Say it:} \emph{suruto, inu ga hoemashita} — At that, the dog began to bark.
  \item \emph{Recall:} say the Japanese for to decrease, then the Japanese for to think of, then say \emph{suruto} again
\end{itemize}

\begin{tcolorbox}[breakable,colback=teal!4,colframe=teal!35!black,title={Before you move on}]
What does \ja{すると} mean? (\textquotedblleft{}Thereupon; at that\textquotedblright{}.) Say it once more.
\end{tcolorbox}
\section[\texorpdfstring{\ja{ところが} (tokoroga)}{ところが (tokoroga)}]{\ja{ところが} (tokoroga) — however; but then}

\label{lesson:JA-C315-tokoroga}



\begin{quote}
Before the new one: say the Japanese for to think of, then the Japanese for thereupon.
\end{quote}

\subsection*{You'll want to know: \ja{ところが}}
\textbf{\ja{ところが}} — \emph{tokoroga} — \textquotedblleft{}however; but then\textquotedblright{}.

\begin{quote}
\textbf{\ja{ところが、へんじは} \ja{きませんでした。}} — \emph{tokoroga, henji wa kimasen deshita} — \textquotedblleft{}But no answer came.\textquotedblright{}
\end{quote}

\begin{cousinweb}
\textbf{\ja{ところが}} (\emph{tokoroga}) is a native Japanese word (wago), inherited rather than borrowed from Chinese: JMdict writes it with the kanji of \emph{tokoro}, \textquotedblleft{}place\textquotedblright{}, and the particle \emph{ga}, and it is usually written in kana (JMdict).
\end{cousinweb}

\begin{grammarlens}[title={What you've built}]
One more word for telling what happened.
\end{grammarlens}

\subsection*{Your turn}
\begin{itemize}
\raggedright
  \item \emph{Say it:} \emph{tokoroga}
  \item \emph{Say it:} \emph{tokoroga}, once more
  \item \emph{Say it:} \emph{tokoroga, henji wa kimasen deshita} — But no answer came.
  \item \emph{Recall:} say the Japanese for to think of, then the Japanese for thereupon, then say \emph{tokoroga} again
\end{itemize}

\begin{tcolorbox}[breakable,colback=teal!4,colframe=teal!35!black,title={Before you move on}]
What does \ja{ところが} mean? (\textquotedblleft{}However; but then\textquotedblright{}.) Say it once more.
\end{tcolorbox}
\section[\texorpdfstring{\ja{しかし} (shikashi)}{しかし (shikashi)}]{\ja{しかし} (shikashi) — but; however}

\label{lesson:JA-C315-shikashi}



\begin{quote}
Before the new one: say the Japanese for thereupon, then the Japanese for however.
\end{quote}

\subsection*{You'll want to know: \ja{しかし}}
\textbf{\ja{しかし}} — \emph{shikashi} — \textquotedblleft{}but; however\textquotedblright{}.

\begin{quote}
\textbf{\ja{しかし、だれも} \ja{いませんでした。}} — \emph{shikashi, daremo imasen deshita} — \textquotedblleft{}But there was nobody there.\textquotedblright{}
\end{quote}

\begin{cousinweb}
\textbf{\ja{しかし}} (\emph{shikashi}) is a native Japanese word (wago), inherited rather than borrowed from Chinese, and it is usually written in kana (JMdict).
\end{cousinweb}

\begin{grammarlens}[title={What you've built}]
One more word for telling what happened.
\end{grammarlens}

\subsection*{Your turn}
\begin{itemize}
\raggedright
  \item \emph{Say it:} \emph{shikashi}
  \item \emph{Say it:} \emph{shikashi}, once more
  \item \emph{Say it:} \emph{shikashi, daremo imasen deshita} — But there was nobody there.
  \item \emph{Recall:} say the Japanese for thereupon, then the Japanese for however, then say \emph{shikashi} again
\end{itemize}

\begin{tcolorbox}[breakable,colback=teal!4,colframe=teal!35!black,title={Before you move on}]
What does \ja{しかし} mean? (\textquotedblleft{}But; however\textquotedblright{}.) Say it once more.
\end{tcolorbox}
\section[\texorpdfstring{\ja{ところどころ} (tokorodokoro)}{ところどころ (tokorodokoro)}]{\ja{ところどころ} (tokorodokoro) — here and there; in places}

\label{lesson:JA-C315-tokorodokoro}



\begin{quote}
Before the new one: say the Japanese for however, then the Japanese for but.
\end{quote}

\subsection*{You'll want to know: \ja{ところどころ}}
\textbf{\ja{ところどころ}} — \emph{tokorodokoro} — \textquotedblleft{}here and there; in places\textquotedblright{}.

\begin{quote}
\textbf{\ja{みちは} \ja{ところどころ} \ja{こおっていました。}} — \emph{michi wa tokorodokoro kōtte imashita} — \textquotedblleft{}The road was icy in places.\textquotedblright{}
\end{quote}

\begin{cousinweb}
\textbf{\ja{ところどころ}} (\emph{tokorodokoro}) is a native Japanese word (wago), inherited rather than borrowed from Chinese, built from \emph{tokoro}, \textquotedblleft{}place\textquotedblright{}, and \emph{dokoro}, \textquotedblleft{}place\textquotedblright{} (from \emph{tokoro}) (JMdict, KANJIDIC2).
\end{cousinweb}

\begin{grammarlens}[title={What you've built}]
One more word for telling what happened.
\end{grammarlens}

\subsection*{Your turn}
\begin{itemize}
\raggedright
  \item \emph{Say it:} \emph{tokorodokoro}
  \item \emph{Say it:} \emph{tokorodokoro}, once more
  \item \emph{Say it:} \emph{michi wa tokorodokoro kōtte imashita} — The road was icy in places.
  \item \emph{Recall:} say the Japanese for however, then the Japanese for but, then say \emph{tokorodokoro} again
\end{itemize}

\begin{tcolorbox}[breakable,colback=teal!4,colframe=teal!35!black,title={Before you move on}]
What does \ja{ところどころ} mean? (\textquotedblleft{}Here and there; in places\textquotedblright{}.) Say it once more.
\end{tcolorbox}
\section[\texorpdfstring{\ja{みはる} (miharu)}{みはる (miharu)}]{\ja{みはる} (miharu) — to keep watch; to guard}

\label{lesson:JA-C315-miharu}



\begin{quote}
Before the new one: say the Japanese for but, then the Japanese for here and there.
\end{quote}

\subsection*{You'll want to know: \ja{みはる}}
\textbf{\ja{みはる}} — \emph{miharu} — \textquotedblleft{}to keep watch; to guard\textquotedblright{}.

\begin{quote}
\textbf{\ja{おまわりさんが} \ja{みちを} \ja{みはっています。}} — \emph{omawarisan ga michi o mihatte imasu} — \textquotedblleft{}A police officer is keeping watch on the road.\textquotedblright{}
\end{quote}

Polite \textbf{\ja{みはります}}, \emph{miharimasu}; a godan verb. Past (polite): \textbf{\ja{おまわりさんが} \ja{えきを} \ja{みはりました。}}, \emph{omawarisan ga eki o miharimashita}.

\begin{cousinweb}
\textbf{\ja{みはる}} (\emph{miharu}) is a native Japanese word (wago), inherited rather than borrowed from Chinese, built from \emph{mi}, \textquotedblleft{}seeing\textquotedblright{} (from \emph{miru}), and \emph{haru}, \textquotedblleft{}to stretch\textquotedblright{} (JMdict, KANJIDIC2).
\end{cousinweb}

\begin{grammarlens}[title={What you've built}]
One more word for telling what happened.
\end{grammarlens}

\subsection*{Your turn}
\begin{itemize}
\raggedright
  \item \emph{Say it:} \emph{miharu}
  \item \emph{Say it:} \emph{miharu}, once more
  \item \emph{Say it:} \emph{omawarisan ga michi o mihatte imasu} — A police officer is keeping watch on the road.
  \item \emph{Recall:} say the Japanese for but, then the Japanese for here and there, then say \emph{miharu} again
\end{itemize}

\begin{tcolorbox}[breakable,colback=teal!4,colframe=teal!35!black,title={Before you move on}]
What does \ja{みはる} mean? (\textquotedblleft{}To keep watch; to guard\textquotedblright{}.) Say it once more.
\end{tcolorbox}
